\section{Ternary Goldbach theorem}

\begin{definition}[Computed candidate triple]
\label{def:tg-triple}
\lean{TernaryGoldbachTheorem.goldbachTriple}
\leanok
A computable candidate triple is used to discharge bounded numerical ranges.
\end{definition}

\begin{lemma}[Verified finite ranges]
\label{lem:tg-small}
\lean{TernaryGoldbachTheorem.smallCases,TernaryGoldbachTheorem.mediumCases}
\leanok
\uses{def:tg-triple}
The theorem is verified computationally for odd integers up to $10000$.
\end{lemma}

\begin{lemma}[Binary Goldbach input]
\label{lem:tg-binary}
\lean{TernaryGoldbachTheorem.evenGoldbachForLarge}
\notready
Every sufficiently large even number is the sum of two primes. This auxiliary statement is a \texttt{sorry}.
\end{lemma}

\begin{lemma}[Large odd integers]
\label{lem:tg-large}
\lean{TernaryGoldbachTheorem.largeCases}
\notready
\uses{lem:tg-binary}
Large odd integers are written as $3$ plus two primes.
\end{lemma}

\begin{theorem}[Ternary Goldbach]
\label{thm:ternary-goldbach}
\lean{TernaryGoldbachTheorem.MainTheorem}
\notready
\uses{lem:tg-small,lem:tg-large}
Every odd integer greater than $5$ is the sum of three primes.
\end{theorem}
